\newpage
\section[Cuckoo Hashing (11)]{Cuckoo Hashing (11) \cite{Pagh2001Cuckoo}} 

\begin{enumerate}
    \item \textit{Category}: Rasmus Pagh and Flemming Friche Rodler’s Cuckoo Hashing is a description of a research prototype. In this paper, they present a simple and new hashing scheme with worst-case constant lookup time, which is very simple to implement, and has average case performance comparable to the best dictionaries implemented at large.

    \item \textit{Context}: To analyze the performance of Cuckoo Hashing in the context of time and space, they primarily studied the theoretical properties of Dietzfelbinger et al.’s dynamic perfect hashing scheme. The introduction of their proposal states that Cuckoo Hashing possesses the same theoretical properties as the classic dictionary of Dietzfelbinger et al., but is much simpler. Apart from simplicity, the paper's abstract reveals that their simple dictionary’s worst-case constant lookup time equals the theoretical performance of the classic dynamic perfect hashing scheme of Dietzfelbinger et al. Thus, besides its other referenced papers, Pagh and Rodler’s paper primarily utilized Dietzfelbinger et al.’s (Dynamic perfect hashing: Upper and lower bounds. SIAM J. Comput., 23(4):738–761, 1994) theoretical dynamic perfect hashing scheme to evaluate Cuckoo Hashing’s general performance.

    \item \textit{Correctness}: The assumptions about Cuckoo Hashing’s performance appear to be valid due to the metrics they considered for performance evaluation. As Dietzfelbinger et al.’s dynamic perfect hashing scheme is a credible basis for analyzing best-performing dictionaries, it follows that after Keshav’s First Pass, the assumptions on Cuckoo Hashing hold.

    \item \textit{Contributions}: The paper’s main contributions are extensive analyses of key arrangement schemes and hashing schemes in ubiquitous computer dictionaries. These analyses were mainly used to evaluate the performance of their principal contribution—Cuckoo Hashing.

    \item \textit{Clarity}: Due to the clear reflection of the paper’s topic in its title and abstract, explicit expression of the paper’s purpose and contribution to its abstract and conclusion, coherence of figures, tables, and references, and finally, cohesiveness of the structured parts of the paper considered so far (title, abstract, info, section, sub-section heading, conclusions, and references), the paper is, thus far, well-written. Indeed, a brief determinant of what made the paper well-written was the ease we experienced in answering these five (5) C’s. 

    
\end{enumerate}

    \textit{Refer to} \hyperref[fig1]{Appendix A1} for a figure on the examples of \texttt{CUCKOO HASHING}

